%% Copernicus Publications Manuscript Preparation Template for LaTeX Submissions
%% ---------------------------------
%% This template should be used for copernicus.cls
%% The class file and some style files are bundled in the Copernicus Latex Package, which can be downloaded from the different journal webpages.
%% For further assistance please contact Copernicus Publications at: production@copernicus.org
%% https://publications.copernicus.org/for_authors/manuscript_preparation.html


%% Please use the following documentclass and journal abbreviations for preprints and final revised papers.

%% 2-column papers and preprints
\documentclass[journal abbreviation, manuscript]{copernicus}



%% Journal abbreviations (please use the same for preprints and final revised papers)


% Advances in Geosciences (adgeo)
% Advances in Radio Science (ars)
% Advances in Science and Research (asr)
% Advances in Statistical Climatology, Meteorology and Oceanography (ascmo)
% Aerosol Research (ar)
% Annales Geophysicae (angeo)
% Archives Animal Breeding (aab)
% Atmospheric Chemistry and Physics (acp)
% Atmospheric Measurement Techniques (amt)
% Biogeosciences (bg)
% Climate of the Past (cp)
% DEUQUA Special Publications (deuquasp)
% Earth Surface Dynamics (esurf)
% Earth System Dynamics (esd)
% Earth System Science Data (essd)
% E&G Quaternary Science Journal (egqsj)
% EGUsphere (egusphere) | This is only for EGUsphere preprints submitted without relation to an EGU journal.
% Earth Observation (eo)
% European Journal of Mineralogy (ejm)
% Geochronology (gchron)
% Geographica Helvetica (gh)
% Geoscience Communication (gc)
% Geoscientific Instrumentation, Methods and Data Systems (gi)
% Geoscientific Model Development (gmd)
% History of Geo- and Space Sciences (hgss)
% Hydrology and Earth System Sciences (hess)
% Journal of Bone and Joint Infection (jbji)
% Journal of Environmentally Compatible Air Transport System (jecats)
% Journal of Micropalaeontology (jm)
% Journal of Sensors and Sensor Systems (jsss)
% Magnetic Resonance (mr)
% Mechanical Sciences (ms)
% Natural Hazards and Earth System Sciences (nhess)
% Nonlinear Processes in Geophysics (npg)
% Ocean Science (os)
% Polarforschung - Journal of the German Society for Polar Research (polf)
% Proceedings of the International Association of Hydrological Sciences (piahs)
% Proceedings of the International Ocean Drilling Programme (piodp)
% Safety of Nuclear Waste Disposal (sand)
% Scientific Drilling (sd)
% SOIL (soil)
% Solid Earth (se)
% State of the Planet (sp)
% The Cryosphere (tc)
% Weather and Climate Dynamics (wcd)
% Web Ecology (we)
% Wind Energy Science (wes)


%% \usepackage commands included in the copernicus.cls:
%\usepackage[german, english]{babel}
%\usepackage{tabularx}
%\usepackage{cancel}
%\usepackage{multirow}
%\usepackage{supertabular}
%\usepackage{algorithmic}
%\usepackage{algorithm}
%\usepackage{amsthm}
%\usepackage{float}
%\usepackage{subfig}
%\usepackage{rotating}

% UNCOMMENT FOR SUBMISSION !!!!!!!!!!!!
\usepackage{xcolor}
% Define the command for comments, one for each author
\newcommand{\bg}[1]{\textcolor{magenta}{#1}}
\newcommand{\ar}[1]{\textcolor{red}{#1}}
\newcommand{\mahe}[1]{\textcolor{teal}{MAHE: #1}}
% ------------------------------------------------

\usepackage[most]{tcolorbox}

\tcbset{
  colback=gray!5,
  colframe=gray,
  boxrule=0.5pt,
  arc=2pt,
  left=6pt,
  right=6pt,
  top=6pt,
  bottom=6pt
}

\newcounter{definitionbox}
\newtcolorbox{definitionbox}[1][]{
  breakable,
  code={\refstepcounter{definitionbox}},
  title=Box \thedefinitionbox: #1
}

\begin{document}

\title{On combining climate models into weighted ensembles }


% \Author[affil]{given_name}{surname}
\Author[1][britta.grusdt@awi.de]{Britta}{Grusdt} %% correspondence author
\Author[1]{Mahé}{Perrette}
\Author[1]{Alexander}{Robinson}

\affil[1]{Alfred Wegener Institute, Helmholtz Centre for Polar and Marine Research, Potsdam, Germany}


%% The [] brackets identify the author with the corresponding affiliation. 1, 2, 3, etc. should be inserted.

%% If an author is deceased, please add \deceased[$Deceased date if applicable$]{$Author number$} (e.g. \deceased[13 November 2015]{2}) at the end of the affiliations. The author number depends on the placement of the author in the author list, e.g. the third author has number 3.


%% If authors contributed equally, please add \equalcontrib{$Author numbers$} (e.g. \equalcontrib{1,3}) at the end of the affiliations. The author number depends on the placement of the author in the author list, e.g. the third author has number 3.




\runningtitle{TEXT}

\runningauthor{TEXT}





\received{}
\pubdiscuss{} %% only important for two-stage journals
\revised{}
\accepted{}
\published{}

%% These dates will be inserted by Copernicus Publications during the typesetting process.


\firstpage{1}

\maketitle



\begin{abstract}
Several methods have been proposed and used to refine estimates of future climate change based on combined output from comprehensive climate models. While previously the so-called model democracy approach was used to combine model predictions, where every model is given equal weight, it is now widely accepted that using model weights that account for model performance and model independence yields better results. However, most existing approaches rely, implicitly or explicitly, on a similar statistical basis, while describing things in different ways. Here we aim to build a common framework formulated in probabilistic Bayesian terms and distinguishing between approaches that are based on the performance of individual models (individual performance weighting) and approaches that are based on the performance of the weighted ensemble as a whole (ensemble performance weighting). Using simple constructed examples, we demonstrate that the ensemble performance weighting approach implicitly accounts for co-dependencies among models, and examine how methodological choices such as the prior distribution affect the results. We attempt to clearly define the characteristics of the task in general terms, and from there, arrive at useful guidelines and strategies to make robust climate projections with uncertainty estimates, as well as establishing a common basis for future developments with regard to ensemble model weighting in Earth system science.
\end{abstract}


\copyrightstatement{TEXT} %% This section is optional and can be used for copyright transfers.


\introduction

\introduction


% 1.Define the task
% What? Ensembles of climate models. Why needed?
For many decades, climate models have been an important tool for gaining a better understanding of the climate system, and for making inferences about future climate and the consequences of climate change. Over the years many different models have become available which is generally beneficial to get a more accurate picture of reality. There are, for instance, different possibilities on how to parameterize processes that cannot be resolved in the models, which leads to uncertainties in the predictions that can be represented by using an ensemble of models. This type of uncertainty is referred to as \emph{model}- or \emph{structural uncertainty}, and it is the primary uncertainty addressed by a multi-model ensemble.
In the most straightforward way, each model in an ensemble is considered equally important and thus receives equal weight. This is therefore known as the ``model democracy'' approach, where simple multi-model means (MMMs) are computed and the ensemble spread is used to quantify uncertainty. While this has long been the standard approach, it comes along with several problems that can be addressed by weighting the models in the ensemble.

% What are the problems with ensembles?
A general challenge is that we must work with ``ensembles of opportunity'': the available models do not represent a statistically perfect ensemble in the sense of a set of unbiased and independent random samples \citep[e.g., see][]{massonClimateModelGenealogy2011, pennellEffectiveNumberClimate2011}. They were not constructed to cover the entire range of plausible outcomes, nor is this space evenly sampled as the models can be dependent on one another in complex and often intransparent ways. Due to such inter-model dependencies, they may, for instance, have shared biases, which essentially means that some models in the ensemble do not provide as much new information as others \citep[e.g.,see][]{junSpatialAnalysisQuantify2012}. Furthermore, the models may differ in how plausible they are in light of the observational data. Therefore, there are two main aspects that weights assigned to models in an ensemble should account for: model performance and model dependence.

% Context: general methods from the literature
Various methods have been proposed in the literature to compute model weights that reflect one or both aspects. 
A method that has received much attention is ClimWIP from \citet{brunnerReducedGlobalWarming2020}, which was developed from a long line of previous work \citep[e.g.,][]{sandersonRepresentativeDemocracyReduce2015, sandersonAddressingInterdependencyMultimodel2015,knuttiClimateModelProjection2017, lorenzProspectsCaveatsWeighting2018}.
It combines weights that are calculated separately for independence, defined in terms of model output similarities, and performance.
Both types of weights are computed based on the assumption of a Gaussian error model and by taking one or more diagnostic variables into account. The diagnostic variables used for computing performance weights may be different from those used to calculate independence weights. 
Meanwhile, the Reliability Ensemble Averaging method (REA) is a similar method introduced by  \citet{giorgiCalculationAverageUncertainty2002a} which also considers model performance and independence.
Another, conceptually different approach is what is often referred to as Bayesian Model Averaging (BMA). While the other methods mentioned above compute model weights based on the performance and/or independence of the individual models in the ensemble, BMA-weights are based on the performance of the weighted ensemble itself. Therefore, the BMA-weights are selected so that the weighted average model simulates the observed data as closely as possible \citep[e,g.][]{massoudBayesianModelAveraging2020,woottenEffectStatisticalDownscaling2020,massoudBayesianWeightingClimate2023,woottenAssessingSensitivitiesClimate2023}. 
It should be noted that this definition of BMA differs from what BMA refers to in the context of statistical models where it was first applied, but since the application of BMA on forecasting surface temperatures by \citet{rafteryUsingBayesianModel2005}, it has become a common way to refer to the method used in at least some climate science studies today.

Although the methods differ in the details, they can mainly be distinguished by whether the ensemble is considered as a whole, which we will refer to as ``\emph{Ensemble Performance Weighting}'', or as individual members, which we will refer to as ``\emph{Individual Performance Weighting}``.
% Questions to be addressed
All approaches to computing model weights have in common that they need to provide answers to the following two questions: (i) how to compute model performance and (ii) how to account for potential dependencies between models.
% Examples for ways to compute performance
Concerning model performance, \citet{sierraInformationtheoreticApproachObtain2025}, for instance, use a method, embedded in information-theoretic concepts, that computes performance weights relative to the inverse of the deviations between model output and data. Similarly, \citet{fengComparisonFourEnsemble2011} use inverse-distance power weighting ($\frac{1}{\textrm{AE}_i^r}, r>0$, AE$_i$: absolute error of model $i$) to compute performance weights. 
Others define performance based on the assumption of independent and identically distributed Gaussian errors \citep[e.g.~][]{brunnerReducedGlobalWarming2020}.
Note that all such individual-performance weighting approaches where the weights are a function of the sum of squared errors yield the same ranking order, only the magnitude of the weights differs (see Appendix \ref{sec:appendix-B} for examples).

% Examples for ways to compute independence
Meanwhile, independence is often defined in terms of model similarities \citep[e.g.~see][]{brunnerReducedGlobalWarming2020}. The more similar the predictions of two models are, the more dependent they are considered to be. This is reasonable in that two dependent models will have shared biases which is reflected in similar model predictions for both. 
Another option is to use meta information about the climate models that could provide further insights about co-dependencies between the models, e.g., information about the modeling group, specific model components or model resolution  \citep[e.g.~see][]{boeInterdependencyMultimodelClimate2018, kumaClimateModelCode2023}. 
However, contrary to the model predictions, this information is not always available, making model similarities a pragmatic choice to define model independence. Furthermore, even if meta information about the models were available, there would not be a unique straightforward definition of independence. As pointed out by \citet{massoudBayesianWeightingClimate2023}, considering full distributions over weight vectors rather than a  single weight per model implicitly provides information about inter-model dependencies. The idea is that this is due to the correlation between the posterior weights of different models. This is an important claim that is fundamental to the distinction between ‘individual’ performance weighting and ‘ensemble’ performance weighting. 

% Short summary paragraph, hinting towards now the main methods sections that come below.
Here we aim to clearly define these two approaches within the common language of probabilities in Bayesian terms, so that we can easily compare and apply both methods in a consistent way. We use simple, idealized examples to test the methods in a systematic way, with a particular focus on the question of how inter-model dependencies are accounted for in a model weighting exercise.  We also  explore how  the results obtained from applying this approach are affected by different methodological choices about the prior distribution over weights, which can be important and should be a prominent aspect of the ensemble performance weighting approach. 

The paper is structured as follows. In Section \ref{sec:framing} we lay out the differences between the two approaches and define the terminology we use. In Sections \ref{sec:priors} and \ref{sec:likelihoods}, we discuss choices for the prior and likelihood distributions and consider the influence of using a set of different diagnostic variables, including a simple example. In Section \ref{sec:dependence} we consider the issue of model dependence through simple examples, while in Section \ref{sec:ecs}, we consider the particular realistic case of using Equilibrium Climate Sensitivity (ECS) as a performance diagnostic. In Section \ref{sec:discussion}, we discuss the overall implications of our work. Section \ref{sec:conclusion} concludes the paper and provides an outlook on future work.






\section{Conceptual framing of weighting approaches}\label{sec:framing}

In the following we will provide definitions for the two general approaches for computing model weights, the individual and the performance weighting approach, and specify the terminology that we use throughout the paper. 
This is summarized in Box \ref{box:terminology}.
Note that generally, lowercase letters refer to predictions/observations for a specific climate variable, while uppercase letters refer to collections, e.g., of different variables or models.
An exception are uppercase calligraphic letters, $\mathcal{M}_i$ and $\mathcal{Y}$, which we use to refer to a model $i$ or an observational dataset as a category.
%Sometimes, we want to refer to different sets of climate variables, which we denote by a set $j$.

\begin{definitionbox}[Terminology]\label{box:terminology}
$\mathcal{M}_i$: model $i$ as an instance/category, e.g.~AWI-ESM 

$m^v_{i}(s)$: predicted variable $v$ for model $i$ at index $s$ (e.g.~spatial location, time point)

$M_i = [m^{\textrm{tas}}_i, m^{\textrm{sst}}_i, \dots]$: set of considered predicted variables for model $i$

$M = [M_1, M_2, \dots]$: superset of predictions from all models

$M_* = \sum_{i} M_i \cdot w_i$ where $\sum_i w_i = 1$: predictions of the weighted average model $\mathcal{M}_*$ 

\bigskip
$\mathcal{Y}$: observational dataset as an instance/category (e.g.~ERA5)

$y^v(s)$: observations for variable $v$ at index $s$

$Y = [y^{\textrm{tas}}, y^{\textrm{sst}}, \dots]$: set of considered observed variables

\end{definitionbox}

\paragraph*{Individual Performance Weighting}

With the individual performance weighting approach, a weight vector is defined such that each weight represents the plausibility of the respective model in light of the observed data $Y$.
In this approach, the weight assigned to a model $\mathcal{M}_i$ is proportional to the product of the likelihood of the observed data $Y$ under model $\mathcal{M}_i$, i.e. given its predictions, $M_i$, and the prior over models, $P(\mathcal{M}_i)$.
%
\begin{align}
    &\label{eq:wi-individual-weighting}  w_i \propto P(Y\mid M_i) \cdot P(\mathcal{M}_i)
\end{align}
%
When each model is \emph{a priori} considered equally likely, the resulting weight vector is only influenced by the likelihood, and the prior can thus be omitted. This simplifies Eq.~\eqref{eq:wi-individual-weighting} into $w_i \propto P(Y\mid M_i)$.
Note the difference between the model as a category, $\mathcal{M}_i$ and its predictions $M_i$, where the former implies the latter: once we decide for a model (e.g.~AWI-ESM), we can look up its predictions. Therefore, for completeness, we note that Eq.~\eqref{eq:wi-individual-weighting} can be spelled out explicitly as 
%
\begin{align}
    &  w_i = P(\mathcal{M}_i \mid Y, M_i) \propto P(Y\mid \mathcal{M}_i, M_i) \cdot P(M_i\mid \mathcal{M}_i)  \cdot P(\mathcal{M}_i) =  P(Y\mid M_i) \cdot 1 \cdot P(\mathcal{M}_i)
\end{align}


\paragraph*{Ensemble Performance Weighting}
In the ensemble performance weighting approach, the objective is to find a weight vector $\mathbf{w}$ so that the likelihood of the observed data is maximized under the weighted average model, $\mathcal{M}_*$.
Thus, performance is not considered for individual models, but for the ensemble as a whole. This is the main difference between both approaches yielding two conceptually different outputs: in the individual weighting approach, we get a discrete distribution over categories (the models in the ensemble), i.e. a single weight vector whereas in the ensemble weighting approach, we get a continuous distribution over weight vectors:
%
\begin{align}
    & \label{eq:w-ensemble-weighting} P(\mathbf{w}\mid Y, M) \propto P(Y\mid M, \mathbf{w}) \cdot P(\mathbf{w}) = P(Y\mid M_*) \cdot P(\mathbf{w})
\end{align}
%
For a weight vector $\mathbf{w}$, the predictions of the weighted average model $\mathcal{M}_*$ for a variable $v$ are given by
%
\begin{align}
    & \label{eq:weighted-avg-model}
    m^v_* = \sum_{i=1}^N w_i \cdot m^v_i
\end{align}
Thus, the log likelihood of the observed data under model $\mathcal{M}_*$, i.e. given its predictions, $M_*$, is computed as
%
\begin{align}
    & \label{eq:logl-m*} \log P(Y\mid M_*) = \sum_v \log P(y\mid m_*^v)
\end{align}

To find the weight vector associated with the maximum likelihood, stochastic sampling methods can be used.
For example, using Markov Chain Monte Carlo (MCMC) sampling, we can approximate the posterior distribution over weights $\mathbf{w}$ (Eq.~\eqref{eq:w-ensemble-weighting}).
This has the advantage that uncertainties around the weights are incorporated and can thus be naturally propagated.
Therefore, instead of assuming a single weight vector when computing predictions (e.g., future climate projections), this approach allows us to compute a posterior predictive distribution that accounts for our uncertainty in the weights.
If we want to summarize the posterior distribution over weights, we may compute the mean or the maximum-a-posteriori (MAP) weight vector.
For unimodal and symmetrical distributions these two values should be nearly identical. 

% Maybe here difference in interpretation of weights for individual vs. ensemble performance weighting?
Note the difference in the interpretation of the weights in the ensemble performance weighting approach compared to the individual performance weighting approach.
The mean estimates of the sampled weights (in the ensemble performance weighting approach) represent the average contribution that a model makes to the weighted average. This is not strictly equivalent to, but can be interpreted as the relative performance of each model.
Contrary to that, the weights computed in the individual performance weighting approach exactly mirror the performance of each individual model compared to the other models in the ensemble.

Irrespective of the applied approach (individual or ensemble performance weighting), we need to define (i) a prior distribution and (ii) a likelihood function (i.e.~$P(Y\mid M_i)$ for each model $\mathcal{M}_i$ or $P(Y\mid M_*$) for the weighted average model, respectively), or more generally speaking, a way to measure model performance with respect to the observed data. We will consider both in turn in the next two sections.

\section{Prior Distributions}\label{sec:priors}
In the individual weighting approach the prior is defined over models (as categories), while in the ensemble weighting approach the prior is defined over weight vectors.
In the absence of any available information about the models before considering the data, like information about model dependencies, for example, the most neutral choice for the prior is a distribution that treats each model, or in the ensemble weighting approach, each weight vector, as equally likely.
In the individual performance weighting approach, this corresponds to an equal probability for each model, or $p_i = \frac{1}{N}$ where $N$ is the number models. Mathematically this is described by a categorical distribution:
\begin{equation}\label{eq:prior-individual} 
    \mathcal{M}_i \sim \textrm{Categorical}(\mathbf{p}=\left[\frac{1}{N},\dots,\frac{1}{N}\right])
\end{equation}
If available, prior knowledge can be taken into account by adjusting the parameter vector $\mathbf{p}$ accordingly, subject to the constraint $\sum_{i=1}^N p_i = 1$.
Analogously,  in the ensemble performance weighting approach, we draw weight vectors from a Dirichlet distribution with parameter vector $\vec{\alpha}$ where we set each entry $\alpha_i$ to the same constant value $c$:
\begin{equation}\label{eq:prior-ensemble}  
    \mathbf{w} \sim \textrm{Dirichlet}(\vec{\alpha} = \left[c,\dots,c\right]) \quad \textrm{with }  c>0.
\end{equation}
%
The Dirichlet distribution is defined over $N$-dimensional vectors that sum up to 1. 
%Since each vector in its support defines a new categorical distribution, it can be considered a distribution over distributions.
%Setting $\alpha_i=1 \forall i$, corresponds to a uniform distribution over all data points (i.e.~vectors) in its support. In this case, 
When all $\alpha_i$ are identical, the vector $\vec{\alpha}$ reduces to a single scalar $\alpha$, also called the concentration parameter.
Figure \ref{fig:priors}a shows the distribution of each component $w_i$ of weight vectors drawn from a Dirichlet distribution over 3-dimensional vectors (corresponding to an ensemble of three models) and how it changes with varying $\vec{\alpha}$. 
The expected value for $w_i$, i.e.~for model $i$, is given by $\frac{\alpha_i}{\alpha_0}$ where $\alpha_0 = \sum_{i=1}^N \alpha_i$. Therefore, when all parameters $\alpha_i$ are set to the same value $c$, the expected value for each $w_i$ is always equal to $\frac{1}{N}$ since $E\left[w_i\right] = \frac{\alpha_i}{\alpha_0} = \frac{c}{N\cdot c} = \frac{1}{N}$. 
This is the case in the two bottom rows of Fig.\ref{fig:priors}a, where $\alpha = 1$ and $\alpha = \frac{1}{N}$ respectively. It is, therefore, a reasonable choice to set all $\alpha_i$ to a single value when we do not have any prior knowledge about the models' weights. Yet, the sampled vectors differ depending on the absolute value of $\alpha$: when $\alpha=1$, the prior mass is evenly spread and all $w_i$ are close to the expected value $\frac{1}{N}$ whereas when $\alpha < 1$, the weights tend to be either large or small (U-shaped distribution). In the asymptotic case where $\alpha \to 0$, the weight vectors approach one-hot vectors, i.e. nearly all probability mass is assigned to a single model, while all others receive almost no weight.
The upper three rows of Fig.~\ref{fig:priors}a further show how the distribution of each $w_i$ changes when an exemplary basis parameter vector of $\left[0.1, 0.4, 0.5\right]$ is multiplied by $0.1, 10$ and $100$ respectively. The larger the parameter values $\alpha_i$ are, the tighter the distributions scatter around the respective mean value.
Therefore, the Dirichlet distribution allows us to easily incorporate potential prior knowledge about the value of a model's weight: the more certain we are \emph{a priori}, the larger we set the value of the corresponding parameter $\alpha_i$. There are alternative distributions for vectors that sum up to 1 such as the Logit Normal distribution. These were, however, not explored in this work.

%An alternative to the Dirichlet distribution is the Logit Normal distribution, which is also defined over vectors that sum up to 1. These are sampled in two steps: first, a vector is drawn from a Multivariate Normal distribution. Then, the softmax function is applied to each sampled vector, which ensures that they sum up to 1. By varying the covariance matrix of the Multivariate Normal distribution, it is, for instance, possible to take into account potential prior knowledge about the correlations between models. 

\begin{figure}
\includegraphics{sections/figures/fig1a.pdf}
\includegraphics{sections/figures/fig1b.pdf}
\caption{(a) Distributions of weights per model sampled from a Dirichlet distribution using different parameters. The stars show the mean weights. (b) Artificial data (grey dots) of 38 ``models'', drawn from $\mathcal{N}(\mu=3,\sigma = 1)$ (black curve). The densities show the distributions of the expected data of the weighted average models computed for each of 10,000 weight vectors drawn from Dirichlet distributions with varying parameters.}
\label{fig:priors}
\end{figure}

Especially in cases where we only have few or not very informative data, the choice of the prior needs careful consideration as, in that case, the posterior predictions will be more influenced by the prior. 
Figure \ref{fig:priors}b shows how the choice of the prior distribution over the weight vectors influences the data that is considered in the first place by the weighted average models. Here, we generated arbitrary data for 38 models (grey stars) by drawing from a Normal Distribution with $\mu=3, \sigma=1$ (black curve) to represent an ensemble of opportunity. 
The three colored densities show the distribution of the weighted averages under Dirichlet priors (over $\mathbf{w}$) with varying concentration parameter $\alpha$. In all three cases, the expected value for each $w_i$ equals $\frac{1}{N}$, reflecting that we do not have any prior knowledge about the models' weights. 
When $\alpha=1$ (green density), all weight vectors that we consider through the prior distribution assign approximately equal probability to all $w_i$, so that we are naturally restricted to weighted averages that sit tightly around the multi-model mean.
When $\alpha$ is extremely small (purple density where $\alpha=\frac{1}{500}$), the considered ``weighted'' averages effectively correspond to the predictions of individual models as the prior distribution generates vectors that approach one-hot vectors (where $\alpha_j=1$, $\alpha_i=0 \quad \forall i\neq j$).
Setting $\alpha = \frac{1}{N}$ (orange density) turns out to be a reasonable choice in that the weight vectors drawn from this prior distribution allow us to cover a large part of the spread of the model ensemble.
%while the expected value of each $w_i$ remains at a reasonable value of $\frac{1}{N}$,  
This is due to the variance of the Dirichlet distribution, with scalar concentration parameter $\alpha$, defined as $var(\mathbf{w}_i)=\lambda \cdot \frac{1}{N \cdot \alpha + 1}$ where $\lambda=\frac{N-1}{N^2}$, i.e.~$\lambda$ does not depend on $\alpha$.
The variance of the prior weights will therefore be very small if $\alpha$ does not depend on the number of models (assuming a moderate number of models, in Fig.~\ref{fig:priors}b, we use $N=38$), even if it is set to a small value like 1.
And since a small variance in the prior weights yields a highly homogeneous set of weight vectors that are considered in the first place, the data that is \emph{a priori} considered by the weighted average models is restricted accordingly.

In summary, in the individual performance weighting approach, a uniform prior over the models will likely be the best choice in most cases, unless one has prior knowledge that justifies to prefer certain models over others. In the ensemble weighting approach, the prior should be considered more carefully as there are several choices with quite different implications even if we assume no prior knowledge, reflected by an expected value for each $w_i$ of $\frac{1}{N}$. In this case, $\alpha = \frac{1}{N}$ seems to be a good choice. 

\section{Measuring performance}\label{sec:likelihoods}

The choice of how to evaluate model performance lies at the heart of every method to compute weights for models in an ensemble.
For data on a spatial grid, an example for a performance metric $f$ is the inverse mean squared error between model predictions and observed data:
%
\begin{equation}
	f(m_i^v,y^v) = \frac{1}{\sum_{s} w_s \cdot (m^v_i(s) - y^v(s))^2}
\end{equation}
%
where $s$ iterates over spatial locations and $w_s$ refers to weights (such as area weighting) normalized to sum to 1. 

In the individual performance weighting approach, the weight for model $i$ is then proportional to the performance of model $i$, based on performance metric $f$, and assuming a uniform prior over models here for simplicity.
%
\begin{equation}
	w_i \propto f(m^v_i, y^v)
\end{equation}
This approach was, for instance, used by \citet{sierraInformationtheoreticApproachObtain2025}. %\bg{add references} 
To formulate the weighting in a Bayesian framework, we need to define a proper generative likelihood function instead of relying on a heuristic performance metric. 
The likelihood relates model predictions to the observed data in a probabilistic way and allows us to incorporate uncertainties that are not accounted for when using a heuristic performance metric.
A possible choice for the likelihood is to use an independent Gaussian error model (i.i.d. errors)
%
\begin{align}
& y^v = m_i^v + \epsilon, \quad \epsilon \sim \mathcal{N}(0, \sigma^2) \\
& y^v \sim \mathcal{N}(m_i^v, \sigma^2)
\end{align}
%
where $m_i$ refers to the output of any model $\mathcal{M}_i$, be it from a weighted-average model or an individual model. Assuming i.i.d. errors is often unrealistic, but is commonly used as a convenient simplification. The performance weights as computed with ClimWIP \citep{brunnerReducedGlobalWarming2020} are an example of assumed Gaussian errors using individual performance weighting, while \citet{massoudBayesianModelAveraging2020}, for instance, use Gaussian errors with ensemble performance weighting. 
Using Gaussian errors requires choosing a value for the variance, $\sigma^2$.
In the individual performance weighting approach, $\sigma^2$ acts as a hyperparameter that regulates the strength of the relative weighting between models: the smaller $\sigma^2$, the fewer models are considered to simulate the observed data well. Thus, the number of models that are assigned weights substantially greater than 0 decreases with decreasing $\sigma^2$. 
%For this reason, the definition of $\sigma^2$ in this context is open to interpretation. 
For ClimWIP, \citet{brunnerReducedGlobalWarming2020} tune a parameter ($\sigma_D$) that essentially represents the variance $\sigma^2$, with the objective to obtain a weighting that is neither too weak (all equal) nor too aggressive (all weight on few models only). As an objective measure of where this value should be, \citet{brunnerReducedGlobalWarming2020} follow \citet{knuttiClimateModelProjection2017} and run perfect model tests with different values for $\sigma_D$. 
They then set $\sigma_D$ to the smallest value such that more than 80\%, (90\% in \citep{knuttiClimateModelProjection2017})  of the perfect models lie in the 10-90\% (5-95\% in \citep{knuttiClimateModelProjection2017}) range predicted by weighting all but the respective ``perfect'' model.
Note, that the ranking of which models are best remains the same irrespective of $\sigma^2$; only the magnitudes of the weights change.
In ClimWIP, the performance weight of a model $i$ is proportional to the likelihood of the observed data under a Gaussian distribution:
\begin{equation}
    \mathbf{w}_i \propto \mathcal{L}(y \mid \mu=m_i^v, \sigma^2 = \frac{\sigma_D^2}{2}) \cdot P(\mathcal{M}_i)
\end{equation}
%
In general, it is also possible to sample $\sigma^2$, which would make the tuning procedure like it was done by  \citet{brunnerReducedGlobalWarming2020}  unnecessary. However, with ClimWIP, or the individual performance weighting approach in general, this would simply result in a nearly equal weighting since $\sigma^2$ would be tuned to increase the likelihood of the observed data under \emph{every} model. Contrary to that, in the ensemble weighting approach, this problem does not arise as the posterior over $\sigma^2$ represents our uncertainty across the weighted-average models, so that there is no direct relation between the individual model weights and $\sigma^2$. A common prior distribution for the unknown variance of a normal distribution is, for instance, the Inverse-gamma distribution, which we also use in the examples in the following sections where we apply the ensemble performance weighting approach. The joint posterior distribution over weight vectors and variance $\sigma^2$ is then given by:
%In the ensemble performance weighting approach, it is not necessary to apply a tuning procedure like  \citet{brunnerReducedGlobalWarming2020} did, as $\sigma^2$ can be reasonably sampled jointly with the weight vectors. We write reasonably, since it is, of course, also possible to sample $\sigma^2$ with ClimWIP or the individual performance weighting approach in general. However, doing so would simply result in a nearly equal weighting since $\sigma^2$ would be tuned to increase the likelihood of the observed data under \emph{every} model. Contrary to that, in the ensemble weighting approach the posterior over $\sigma^2$ represents our uncertainty across all weighted-average models. 
%The larger $\sigma^2$ is, the more the predictions of the weighted averages deviate from the observations, while smaller values for $\sigma^2$ indicate that the predictions of the weighted average models are closer to the observations.%\footnote{In Sec.~\ref{sec:appendix-sigma}, we go in a bit more detail consider the influence of $\sigma^2$ and how it is related to the weights $\mathbf{w}$.}
%
\begin{equation}
    P(\mathbf{w}, \sigma^2\mid M) \propto \mathcal{L}(y \mid \mu=m_*^v, \sigma^2) \cdot P(\mathbf{w}) \cdot P(\sigma^2)
\end{equation}
where $m_*^v$ are the predictions for variable $v$ of the weighted average model using weight vector $\vec{w}$.
%
Setting $\sigma^2$ to a fixed value in the ensemble weighting approach is also an option.  \citet{massoudBayesianModelAveraging2020}, for instance, use the following log likelihood function: $-\frac{1}{2} \sum_{s} (Y(s) - M_i(s))^2$, where $s$ iterates over all grid cells\footnote{\citet{massoudBayesianModelAveraging2020} refer to this function (their Eq.~(4)) as the likelihood, but it seems that the log likelihood is intended.}, which is equivalent to assuming Normal Gaussian errors centered around 0 with variance $\sigma^2=1$.


\paragraph*{Measuring performance based on a set of diagnostics}\label{sec:measuring-performance-combined-diagnostics}

In the previous section, we focused on a single diagnostic $v$. This could, for instance, be the temperature anomaly of a specific time period. 
In this section, we want to consider the case when the performance weights are based on more than one diagnostic.
There are generally two approaches: either the weights are computed separately for each diagnostic and are then averaged to yield a final weight vector or the different diagnostics are considered all at once.
In the ClimWIP-method (individual weighting) for example, \citet{brunnerReducedGlobalWarming2020} compute one distance value for every model across a set of five different diagnostics. To make the diagnostics unitless, that is, to bring them all on the same scale, they normalize the distance value for each model by the median across all models, respectively for every diagnostic. 
The normalized distances are summed up for every model so that eventually each model is assigned one ``generalized'' distance value.
Performance weights are then based on a Gaussian distribution over these values. In this way, the diagnostics are considered together. An alternative example for using several diagnostics, yet in the  ensemble weighting approach, is the study from \citet{massoudBayesianModelAveraging2020}, who compute one weight vector separately for each diagnostic and then take the average across all computed weight vectors.

To show the implications of this choice --- i.e.~whether to compute one distribution of weights for every diagnostic or a single distribution of weights that takes into account all diagnostics at once --- we generated toy data for two models and two diagnostics, shown in Fig.~\ref{fig:toy-example-diagnostics}.
The data that we use here are anomalies of climatologies for the time period from 1980--2014 with respect to the global mean value. For diagnostic 1, we use near surface air temperature (tas) and for diagnostic 2, we use sea level pressure (psl). To bring the data of both diagnostics on the same scale, we standardize them by the standard deviation of the observations (ERA5) for the respective diagnostic.
The toy model data were constructed so that for diagnostic 1, $\mathcal{M}_1$ is much better than $\mathcal{M}_2$, while for diagnostic 2, $\mathcal{M}_2$ is better than $\mathcal{M}_1$, albeit only marginally (see Fig.~\ref{fig:toy-example-diagnostics}).
In this way, $\mathcal{M}_1$ is objectively the better model and thus, when both diagnostics are taken into account at the same time, we expect $\mathcal{M}_1$ to receive larger weights on average than $\mathcal{M}_2$.
Nonetheless, as $\mathcal{M}_2$ performs better for diagnostic 2 than $\mathcal{M}_1$, it should still payoff to include $\mathcal{M}_2$ to some extent in a weighted average than relying only on $\mathcal{M}_1$ alone.

To compute the posterior distributions over weight vectors, we use a Dirichlet prior parameterized with $\vec{\alpha} = [\frac{1}{2}, \frac{1}{2}]$ and a multivariate Gaussian likelihood with the model predictions as mean vector and variances $\sigma_{d_1}^2, \sigma_{d_2}^2$ for diagnostic 1 and 2 respectively that we sample jointly with the weights. Note that the values for the variances are not reported further, as they are not relevant to the result. 
The results summarized in Table \ref{tab:toy-example-diagnostics-rmses} confirm our expectations: when both diagnostics are jointly used, the mean weight vector is $\left<\bar{w}_1=0.79, \bar{w}_2=0.21\right>$, giving more weight to model $\mathcal{M}_1$ than to model $\mathcal{M}_2$. 
Also, the summed RMSE for the corresponding weighted average model is smaller (0.21) than it is when using the first model alone (0.24) because of the reduced RMSE for diagnostic 2 (0.11 vs. 0.14) that results from giving some weight also to model $\mathcal{M}_2$.

\begin{figure}
    \centering
    \includegraphics{sections/figures/fig2.pdf}
    \caption{Generated toy data for two models and diagnostics so that model 1 is the objectively better model. All data were normalized by the standard deviation of the corresponding observations. The grey lines show the 1:1 perfect matches. (a) Diagnostic 1 with $m_1^{\textrm{tas}} = y^{\textrm{tas}} + \mathcal{N}(0, s_{\textrm{tas}})$ and $m_2^{\textrm{tas}} = y^{\textrm{tas}} + \mathcal{N}(0, 3\cdot s_{\textrm{tas}})$.
    (b) Diagnostic 2 with $m_1^{\textrm{psl}} = y^{\textrm{psl}} + \mathcal{N}(0, 1.5\cdot s_{\textrm{psl}})$ and $m_2^{\textrm{psl}} = y^{\textrm{psl}} + \mathcal{N}(0, s_{\textrm{psl}})$. The standard deviations for the noise were set to $s_{\textrm{tas}}=2, s_{\textrm{psl}}=90$.
    }
    \label{fig:toy-example-diagnostics}
\end{figure}

When instead, both diagnostics are considered separately, i.e.~we first compute the posterior distributions over weight vectors, one based on diagnostic 1, the other based on diagnostic 2, and then average the resulting mean posterior weights, we get a mean weight vector of $\left<\bar{w}_1=0.6, \bar{w}_2=0.4\right>$, thus reducing the weight assigned to model $\mathcal{M}_1$ compared to the joint case.
With this weight vector, the summed RMSE is larger (0.23) than it is when using both diagnostics jointly (0.21), yet still smaller than when using only model $\mathcal{M}_1$ (0.24).

This example shows that, when possible, it is better to consider all diagnostics simultaneously to obtain more representative weight vectors.

\begin{table}[]
    \centering
    %\begin{tabular}{lcccccccc}
     %    diagnostic variables for BMA & $\bar{\sigma}_{\textrm{tas}}$ & $\bar{\sigma}_{\textrm{psl}}$ & $\bar{w}_{\mathcal{M}_1}$ & $\bar{w}_{\mathcal{M}_2}$ & RMSE($m^{\textrm{tas}}_*, y^{\textrm{tas}}$) & RMSE($m^{\textrm{psl}}_*, y^{\textrm{psl}}$) & $\sum$ RMSE \\\hline\hline
     %   tas  & 0.11 & - & 0.90 & 0.10 & 0.10 & (0.13) & 0.23\\
    %    psl  & - & 0.10 & 0.30 & 0.70 & (0.21) & 0.08 & 0.29 \\ \hline
    %    tas + psl (jointly)  & 0.11 & 0.12 & 0.78 & 0.22 & 0.10 & 0.11 & $\mathbf{0.21}$\\
     %   tas, psl (separately)  & 0.97 & 0.84 & 0.60  & 0.40 & 0.14 & 0.09 & 0.23 \\\hline
     %   -  & - & - & 0.5  & 0.5 & 0.16 & 0.09 & 0.25 \\
     %   -  & - & - & 1  & 0 & 0.10 & 0.14 & 0.24 \\
     %   -  & - & - & 0  & 1 & 0.30 & 0.09 & 0.39
      %  \\\hline\hline
      %  \\
    %\end{tabular}
    %\begin{tabular}{lcccccccc}
     %    diagnostic variables for BMA & $\bar{\sigma}_{\textrm{tas}}$ & $\bar{\sigma}_{\textrm{psl}}$ & $\bar{w}_{\mathcal{M}_1}$ & $\bar{w}_{\mathcal{M}_2}$ & RMSE($m^{\textrm{tas}}_*, y^{\textrm{tas}}$) & RMSE($m^{\textrm{psl}}_*, y^{\textrm{psl}}$) & $\sum$ RMSE \\\hline\hline
      %  tas  & 0.05 & - & 0.89 & 0.11 & 0.10 & (0.13) & 0.23\\
      %  psl  & - & 0.05 & 0.32 & 0.68 & (0.21) & 0.08 & 0.29 \\ \hline
       % tas + psl (jointly)  & 0.05 & 0.05 & 0.78 & 0.22 & 0.10 & 0.11 & $\mathbf{0.21}$\\
        %tas, psl (separately)  & 0.05 & 0.05 & 0.60  & 0.40 & 0.14 & 0.09 & 0.23 \\\hline
        %-  & - & - & 0.5  & 0.5 & 0.16 & 0.09 & 0.25 \\
        %-  & - & - & 1  & 0 & 0.10 & 0.14 & 0.24 \\
        %-  & - & - & 0  & 1 & 0.30 & 0.09 & 0.39
        %\\\hline\hline
        %\\
    %\end{tabular}
        \begin{tabular}{lcccccc}
        \tophline
         diagnostic variables & $\bar{w}_1$ & $\bar{w}_2$ & RMSE($m^{\textrm{tas}}_*, y^{\textrm{tas}}$) & RMSE($m^{\textrm{psl}}_*, y^{\textrm{psl}}$) & $\sum$ RMSE\\ \middlehline
        tas  & 0.88 & 0.12 & 0.10 & (0.13) & 0.23\\
        psl  &  0.32 & 0.68 & (0.21) & 0.08 & 0.29 \\ \middlehline
        tas + psl (jointly)  & 0.79 & 0.21 & 0.10 & 0.11 & $\mathbf{0.21}$\\
        tas, psl (separately)  & 0.60  & 0.40 & 0.14 & 0.09 & 0.23 \\ \middlehline
        -  &  0.5  & 0.5 & 0.16 & 0.09 & 0.25 \\
        -  &  1  & 0 & 0.10 & 0.14 & 0.24 \\
        -  &  0  & 1 & 0.30 & 0.09 & 0.39 \\ \bottomhline
    \end{tabular}
    \belowtable{} 
    \caption{Mean posterior weights (columns 2--3, values are rounded) with RMSEs for the corresponding weighted average model, using the constructed data from Fig.~\ref{fig:toy-example-diagnostics}. The RMSEs for the diagnostics that were not used in the objective function are put in parentheses (columns 4--5). Note the weight vector in row 4 (tas,psl separately) corresponds to $\frac{w^{\textrm{tas}} + w^{\textrm{psl}}}{2}$, i.e. it is not sampled, but the average of weight vectors in rows 1--2, where each diagnostic was considered independently. For comparison, row 5 shows the RMSEs for the multi-model mean and the last two rows show the RMSEs when relying on individual models, using either model 1 or model 2.}
    \label{tab:toy-example-diagnostics-rmses}
\end{table}



\section{Model dependence}\label{sec:dependence}

We would like to use three schematic (toy) examples to illustrate how the possible lack of independence between ensemble members can influence the weights obtained when using individual or ensemble performance weighting. 
In each case, we use an ensemble of two base models, $\mathcal{M}_1$ and $\mathcal{M}_2$, whose predictions are systematically constructed based on the 1980--2014 mean near-surface air temperature anomalies with respect to the global mean, obtained from ERA5 \citep{hersbachERA5GlobalReanalysis2020a}. %for the climatology for the same time period.

In the first case, we assume that the first model ($\mathcal{M}_1$) makes perfect predictions in the Eastern Hemisphere, while in the Western Hemisphere, a random error sampled from $\mathcal{N}(0,2)$ is added to the field. For the second model ($\mathcal{M}_2$), the same random error is added in the Eastern Hemisphere, and perfect predictions are obtained in the West. In this way, we get two models with the exact same mean squared error (MSE) with respect to the observed data and we know that the weighting that would minimize the total combined error would be $w_1=w_2=0.5$. The data for this case are shown in Fig.~\ref{fig:independence-error-pattern}.\footnote{For better visualization, Fig.~\ref{fig:independence-error-pattern}--\ref{fig:independence-real-data} exclude the polar regions, but all data was used in the computations.}

\begin{figure}
    \centering
    %\includegraphics{sections/figures/independence/data-error-pattern.pdf}
    \includegraphics{sections/figures/fig3.pdf}
    \caption{Generated toy data with same error pattern for both models. (a) Observations, near-surface air temperature anomalies in the time between 1980 and 2014 with respect to the global mean value of this time period. (b) $M_1(s) = Y(s) + \mathcal{N}(0,2)$ where $s$ is a grid point in the Western hemisphere. For all $s$ in the Eastern hemisphere, the predictions of model 1 are identical to the observations. (c) Predictions of model 2, with identical errors as for model 1, yet in the Eastern hemisphere, while predictions in the Western hemisphere are perfect.}
    \label{fig:independence-error-pattern}
\end{figure}

In the second case, we set the predictions for the two models everywhere to the observed data and independently add the noise to both models (again sampled from $\mathcal{N}(0,2)$). Therefore, in this example, the mean squared errors will be similar, but are very unlikely to cancel out at every grid point. The constructed data for this case are shown in Fig.~\ref{fig:independence-random-error}.

\begin{figure}
    \centering
    %\includegraphics{sections/figures/independence/data-random-err.pdf}
    \includegraphics{sections/figures/fig4.pdf}
    \caption{Generated toy data with independent noise per model. (a) Observations, near-surface air temperature anomalies in the time between 1980 and 2014 with respect to the global mean value of this time period. (b) $M_1(s) = Y(s) + \mathcal{N}(0,2)$. (c) Predictions of model 2, analogous to model $\mathcal{M}_1$, but with independently sampled noise.}
    \label{fig:independence-random-error}
\end{figure}

In the third example, we do not construct data, but use the predictions for the same variable and diagnostic from two real models, namely from AWI-CM-1-1-MR and CESM2-WACCM. These data are shown in Fig.~\ref{fig:independence-real-data}.

\begin{figure}
    \centering
    %\includegraphics{sections/figures/independence/data-real-models.pdf}
    \includegraphics{sections/figures/fig4.pdf}
    \caption{Observations of near-surface air temperature anomalies in the time between 1980 and 2014 with respect to the global mean value of this time period (a) and corresponding predictions of two real models, AWI-CM-1-1-MR (b) and CESM2-WACCM (c).}
    \label{fig:independence-real-data}
\end{figure}

In each of the three toy examples described above, we build ensembles of different sizes by adding copies of model $\mathcal{M}_2$. We know that, by construction, the correct weight of model $\mathcal{M}_1$ (included just once) will always have the same value, irrespective of the number of copies of model $\mathcal{M}_2$ included in the ensemble (and conversely, the sum of weights attributed to the included copies of model $\mathcal{M}_2$ should be equal to the constant value $1-w_1$). In the first toy example, both models have the exact same performance (same MSE) so that the weight for $\mathcal{M}_1$ should remain at $\frac{1}{2}$, while the individual copies of model $\mathcal{M}_2$ should each receive a weight of $\frac{1}{2}\cdot \frac{1}{n_2}$ where $n_2$ is the number of included copies for $\mathcal{M}_2$. The same pattern should apply to the other two toy examples, but with the weight for $\mathcal{M}_1$ remaining at whatever value it receives when model $\mathcal{M}_2$ is included in the ensemble only once. In the second toy example, where we use random errors, this should still be close to $\frac{1}{2}$ whereas in the third example with real model data it depends on how much better or worse $\mathcal{M}_1$ is compared to $\mathcal{M}_2$. In the case that the model weights deviate from these expectations, it implies that the known (strong) dependence between the copies of model $\mathcal{M}_2$ is not properly accounted for.

\begin{figure}
    \centering
    %\includegraphics{sections/figures/independence/independence-copies-mse-gaussian-dirichlet-1-over-n.pdf}
    \includegraphics{sections/figures/fig6.pdf}
    \caption{Model weights for models in an ensemble consisting of model $\mathcal{M}_1$ and $n_2$ copies of model $\mathcal{M}_2$. Panels a-c show the computed weights in the individual performance weighting approach, panels d-f show the mean posterior weights in the ensemble performance weighting approach. Columns refer to the three different toy examples. Toy example 1 (a+d): $\mathcal{M}_1$ and $\mathcal{M}_2$ predict one half of the data perfectly and show the same bias ($\mathcal{N}(0,2)$) in the other half. Toy example 2 (b+e): Model predictions are constructed from the observed data plus some random noise everywhere: $m_i = y + \epsilon$ where $\epsilon \sim \mathcal{N}(0, 2)$. Toy example 3 (c+f): $\mathcal{M}_1$ and $\mathcal{M}_2$ correspond to AWI-CM-1-1-MR and CESM2-WACCM. The two values in brackets in the titles are the mean squared errors for models $\mathcal{M}_1$ and $\mathcal{M}_2$ in the respective example.}
    \label{fig:toy-examples-results}
\end{figure}

Figure \ref{fig:toy-examples-results} shows the resulting weights for these toy examples, using individual or ensemble performance weighting. In the individual performance weighting approach, we evaluate the models by their MSEs (identical to using the sum of squared errors). In the ensemble performance weighting approach, we use a Gaussian likelihood function with the variance jointly sampled with the weight vectors; the shown weights are the mean posterior weights. For the prior over weight vectors, we use a Dirichlet distribution with $\alpha_i = \frac{1}{N} \forall i$ where $N$ is the respective total number of models in the considered ensemble. 
%In Section \ref{sec:ecs}, we look at the influence of the prior, and particularly show why this parameterization should be preferred over using $\mathbf{1}$, which might seem to be a natural choice too.

\begin{figure}
    \centering
    %\includegraphics{sections/figures/independence/independence-copies-posterior-weights-gaussian-dirichlet-1-over-n.pdf}\qquad
    %\includegraphics{sections/figures/independence/independence-copies-correlations-dirichlet-1-over-n.pdf}
    \includegraphics{sections/figures/fig7a.pdf}\qquad
    \includegraphics{sections/figures/fig7b.pdf}
   \caption{Posterior samples using the ensemble performance weighting approach for a toy ensemble of 3 models and with Gaussian errors centered around 0; the variance is jointly sampled with the weight vectors. Green stars show the mean values.}
    \label{fig:toy-examples-posterior-weights}
\end{figure}

In the individual weighting approach, the weights for model $\mathcal{M}_1$ decrease as $N$ increases, essentially following the relationship $w_1=\frac{1}{N}$.
In this case the weight given to any individual copy of model $\mathcal{M}_2$ follows the same relationship. This means that the total weight assigned to  $\mathcal{M}_2$ increases to $\frac{N-1}{N}$.
%That is, as the number of copies of model $\mathcal{M}_2$ increases, model $\mathcal{M}_2$ is given more weight overall in the ensemble compared to model $\mathcal{M}_1$. 
Conversely, in the ensemble weighting approach, neither the weight for model $\mathcal{M}_1$ nor the summed weight across all copies of model $\mathcal{M}_2$ changes with increasing $N$.
%\footnote{When using a Dirichlet prior with all $\alpha_i=1$, we do see a bit of a drop of the weight assigned to model $\mathcal{M}_1$ when as much as 10 copies of model $\mathcal{M}_2$ are included in the ensemble using real data (lower right plot, Fig.~\ref{fig:toy-examples-results}).
%This can, however, be explained by the prior that becomes more influential when more copies are included. With a larger number of copies, increase the number of parameters (weights) to be inferred, but the additional data (the copies of $\mathcal{M}_2$) does not provide any new information.
%Given that the decrease of the weight for model $\mathcal{M}_1$ is marginal compared to the decrease observed in the individual performance weighting approach and given the rather artificial assumption of 10 exact copies of one model, it is still justifiable to conclude that the ensemble weighting approach is able to account for the dependencies among the models, which we introduced by adding $n$ copies of model $\mathcal{M}_2$ to the ensemble.} 
Thus, only the ensemble weighting approach is able to account for the dependencies among the models in the ensemble that we introduced by adding copies of model $\mathcal{M}_2$ to the ensemble.
When performance is measured for the entire ensemble, adding copies of one model does not have an effect on the mean weights since no new information is added. In the individual weighting approach, a newly added copy, however, is seen to add information to the ensemble by definition as every model is considered separately.
Implicitly accounting for model dependence is a strong advantage of the ensemble weighting approach, making the additional computation of independence weights unnecessary.
Nonetheless, in the individual weighting approach, it is possible to combine performance and separately determined independence weights so that the weight of dependent models is reduced \citep[e.g.~see tests from][]{brunnerReducedGlobalWarming2020}.

% Interpretation of model weights
To show why the ensemble weighting approach implicitly accounts for model dependency, Fig.~\ref{fig:toy-examples-posterior-weights} shows the posterior weights for our third toy example, using real data of two CMIP6 models in an ensemble with three members, including $\mathcal{M}_1$ once and $\mathcal{M}_2$ twice. 
We can see that there is a strong correlation between the weights of the two copies of model $\mathcal{M}_2$ (Fig.~\ref{fig:toy-examples-posterior-weights}a).
If the weight of the first copy is high, then the weight of the second copy is necessarily low. 
This effect was also shown for a simpler case by \citet{massoudBayesianModelAveraging2020}.
Contrary to that, the weight of (each copy of) model $\mathcal{M}_2$ does not depend on the weight of model $\mathcal{M}_1$ except for the constraint that the weights of all models have to sum up to 1.
This is also reflected in the variance of the posterior weights (Fig.~\ref{fig:toy-examples-posterior-weights}b): while the variance for the weights of the two copies of model $\mathcal{M}_2$ is large, the variance for the independent model ($\mathcal{M}_1$) is quite small.

Furthermore, remember that any single weight vector sampled from the posterior distribution does not reflect each model's performance like it does in the individual performance weighting approach. For example, when $w_1 = 0.8, w_{2.1}=0.2$ and $w_{2.2}=0.0$ where $w_{2.1}$ and $w_{2.2}$ refer to the weights for the two copies of model $\mathcal{M}_2$, this does not mean that the first copy of model $\mathcal{M}_2$ shows a better performance than the second copy of the same model with respect to the observed data. Nonetheless, the average weight vector does reflect the relative performance of each model within the given ensemble. In this example, the average weight of each copy of $\mathcal{M}_2$ is the same, reflecting that each copy of this model shows equivalent performance (or contributes the same information to the ensemble mean). This difference between sampled weight vectors and the mean weight vector highlights how the conceptualization of the problem is fundamentally different than when individual performance weighting is used.


\section{Performance based on Equilibrium Climate Sensitivity}\label{sec:ecs}

ECS (Equilibrium Climate Sensitivity) as well as TCR (Transient Climate Response) are promising metrics to compute model weights as they are strongly related to future climate change and most of the uncertainty in future climate projections is attributable to the uncertainty in ECS \citep{sherwoodAssessmentEarthsClimate2020,armourEarthsEnergyBudget2021}.
Here, we will focus on ECS and show how to use it as performance metric in a fully Bayesian way. 
ECS has been used previously to compute model weights, for example by \citet{massoudBayesianWeightingClimate2023} who also used an ensemble weighting approach, as we do here, yet without spelling it out completely in a Bayesian framework.

We aim to infer the posterior probability of weight vectors $\mathbf{w}$ given model predictions $M$ (here referring to the models' ECS values):
%
\begin{equation}
    P(\mathbf{w}\mid M) \propto P(M^*\mid \mathbf{w}) \cdot P(\mathbf{w})
\end{equation}
%
We use all CMIP6 models for which we have ECS values, as well as near-surface-air temperatures for SSP585 projections and the historical time period from 1850-1900 available ($N=38$). % the data needed to approximate a model's ECS value with the Gregory method (piControl, 4xCO2 for near-surface-air temperature) as well as projection data for near-surface-air temperatures.
As likelihood function, i.e.~to evaluate the ``performance'' of a weighted average model, that we also refer to as $\mathcal{M}_*$ model, we use the distribution of likely ECS values that \citet{sherwoodAssessmentEarthsClimate2020} derived based on various lines of evidence.
For the prior distribution over weight vectors $\mathbf{w}$ we use a Dirichlet distribution. As we showed in Section \ref{sec:priors}, the choice of the parameterization of the Dirichlet prior distribution is important since it influences the range of the considered data (here ECS values) of the resulting $\mathcal{M}_*$ models.
Ideally, we want the prior weights to span the entire range between 0 and 1 so that a large part of the individual models' ECS values is covered by the prior weighted average models. We achieve this by setting the concentration parameter $\alpha$ to $\frac{1}{N}$. Figure \ref{fig:expected-ecs-dirichlet} shows the distribution of the ECS values for the weighted average models ($\mathcal{M}_*$), based on the prior (pink density) and posterior weights (green density), using Dirichlet($\alpha=\frac{1}{N}$) (left) and, for comparison, using Dirichlet($\alpha=1$) (right). 
Only when the prior is sufficiently broad (Fig.~\ref{fig:expected-ecs-dirichlet}a), does the mean ECS value of the $\mathcal{M}_*$ models, computed based on the posterior weights, shift towards the likely range of the target distribution at around $3^\circ$ Celsius, below the multi-model mean.
When $\alpha = 1$ (Fig.~\ref{fig:expected-ecs-dirichlet}b), the range of $\emph{a priori}$ considered ECS values tightly sits around the multi-model mean as, with this parameterization, on average each model receives equal weight. 


%When all $\alpha_i=\frac{1}{N}$, the weight vectors tend to be more sparse such that most probability mass is assigned to just a few models, which in turn, broadens the range of \emph{a priori} considered ECS values.

% I think since we have the other plot with the prior distributions in the beginning now, we can skip the extra prior plots here
\begin{figure}
    \centering
    %\includegraphics{sections/figures/ecs/expected-ecs-posterior-dirichlet-distributions.pdf}
    \includegraphics{sections/figures/fig8.pdf}
    \caption{Distributions of expected ECS values of weighted ensembles with $N=38$ CMIP6 models, based on prior (pink) and posterior (green) weights, using a Dirichlet prior distribution with (a) $\alpha = \frac{1}{N}$ and with (b) $\alpha = 1$ (right). The black curve is the likelihood function and represents the estimated range of ECS values from \citet{sherwoodAssessmentEarthsClimate2020}.
    }
    \label{fig:expected-ecs-dirichlet}
\end{figure}


\begin{figure}
    \centering
    %\includegraphics{sections/figures/ecs/projections-bma-mmm-with-prior.pdf}
    \includegraphics{sections/figures/fig9.pdf}
    \caption{Predicted increase in global mean near-surface air temperature with respect to the time period from 1995--2014 until the end of the century, under emission scenario SSP5.85, using an ensemble of $N=38$ CMIP6 models. Thick lines show the observations (black), the multi-model mean (orange) and the weighted averages, using the mean posterior (blue) and mean prior (dashed purple) weights when applying the ensemble performance weighting approach with prior Dirichlet($\mathbf{\frac{1}{N}}$) and with the target distribution from Fig.~\ref{fig:expected-ecs-dirichlet} as likelihood. Uncertainty bands (shaded regions) represent the 0.05 and 0.95 quantiles of the predicted (weighted) ensemble spread.}
    \label{fig:projections-mmm-bma}
\end{figure}

We use the computed posterior over weight vectors and apply it to future projections of the increase in global mean near-surface air temperature. 
%\begin{equation}
 %   P(y^{\Delta\textrm{tas}}\mid M^{\Delta\textrm{tas}}) = \int_{\mathbf{w}} P(M^{\Delta\textrm{tas}}_* \mid M^{\Delta\textrm{tas}}, \mathbf{w}) \cdot P(\mathbf{w}) d\mathbf{w}
%\end{equation}
%
%where $\Delta {\textrm{tas}}$ refers to the anomalies with respect to the time period from 1995--2014.
As shown in Fig.~\ref{fig:projections-mmm-bma}, the expected increase in global mean temperatures is lower in the weighted (blue) than in the unweighted (orange) ensemble which we would hope to see, given the known tendency of CMIP6 models to overestimate future warming \citep[e.g.~][]{tokarskaWarmingTrendConstrains2020,hausfatherClimateSimulationsRecognize2022}.
The uncertainty range (inner 90\% quantile) is substantially reduced in the weighted ensemble as compared to the unweighted ensemble spread.
A reduced uncertainty range is expected by the mere fact that we consider weighted averages. 
For comparison, the thin dashed green lines in Fig.~\ref{fig:projections-mmm-bma}, show the 0.05 and 0.95 quantiles for the prior weighted ensemble, i.e.~using the prior weights (Dirichlet($\alpha = \frac{1}{N}$)). The 90\% inner quantile range of the unweighted ensemble is in between 1.26 and 1.56 times larger than the 90\% inner quantile range of the prior weighted ensemble (mean 1.36). 
This shows how much of the reduced uncertainty of the posterior weighted ensemble is just due to the choice of the prior distribution, and how much can be attributed to the data.
In any case, it is important to note that the weighting is not necessarily better the smaller the resulting uncertainty range is.
If, in the ensemble weighting approach, the prior strongly restricts the considered weighted averages (as in Fig.~\ref{fig:expected-ecs-dirichlet}b), this will influence the uncertainty range of the weighted ensemble for the target variable in that it will be unreasonably small. Therefore, in the ensemble performance weighting approach, prior predictive checks become indispensable for a reasonable interpretation of the weighted ensemble. Figure \ref{fig:appendix-projections-mmm-bma} in the appendix (sec.~\ref{sec:appendix}) shows the drastically reduced uncertainty range for the weighted projections when using a Dirichlet prior with $\alpha=1$, i.e.~corresponding to the data in Fig.~\ref{fig:expected-ecs-dirichlet}b. 

How skillful a weighting approach generally is compared to using an unweighted ensemble should furthermore always be assessed through tests \citep[][]{abramowitzESDReviewsModel2019}. If the computed weights are based on historical data, perfect model tests are an option, where one model is removed from the ensemble, and its predictions are considered as ``observations''. It can then be checked whether the pseudo-observations lie in the predicted range for future timesteps. This procedure is then repeatedly done, each time using a different model as ``perfect model'' \citep[e.g.,see][]{brunnerReducedGlobalWarming2020}. Another option is classical calibration-validation where a subset of the observed data is left out for computing the weights. These data can then be used to test whether the weighted ensemble predicts the out-of-sample data better than the unweighted ensemble. Using ECS as performance metric, as we do here, is different in that it involves an independently derived likelihood function. 
In this case, one option is to check whether the ECS-based weighting is valuable for weighting the model predictions for the historical period.
As expected after consideration of the results in Fig.~\ref{fig:projections-mmm-bma}, the ECS-based weighting is not particularly valuable to weigh the historical period.  The summed RMSEs between the observations and the predicted mean values across all years between 1950 and 2014, is nearly identical for the unweighted and the weighted $M_*$ model using the mean posterior weight vector. With the weighted ensemble, the RMSE is only marginally smaller (37.68 vs. 37.83 or on average per timestep 0.5796 vs.0.5820). For future projections, ECS will however become more important as we see in Fig.~\ref{fig:projections-mmm-bma}, where the ECS-weighting constrains the predictions more in the future than it does in the historical period.


% Estimation of ECS based on Gregory method:
\begin{comment}
Climate sensitivity measures the size of the global average surface air temperature response to a radiative forcing applied to the climate system \citep{gregoryNewMethodDiagnosing2004}.
In its original definition, Equilibrium Climate Sensitivity (ECS) refers to the global mean surface air temperature change after an instantaneous doubling of the atmospheric CO$_2$ concentration after reaching equilibrium.
As this is computationally expensive to compute, ECS is usually approximated by the ``effective climate sensitivity'' \citep{schlundEmergentConstraintsEquilibrium2020} using the Gregory method \citep{gregoryNewMethodDiagnosing2004}. 
This is computed based on the first 150 years after an instantaneous quadrupling of the atmospheric CO$_2$ concentration, namely by regressing the top-of-atmosphere net downward radiation $N$ against the global mean surface air temperature change $\Delta T$. In a steady-state climate, $N$ equals 0.
%
\begin{align}
    & N = F - H \\
    & \label{eq:ecs-gregory} N \approx F - \alpha \Delta T
\end{align}
%
where $F$ is the effective radiative forcing from a quadrupling of the atmospheric CO$_2$ concentration, $H$ is the radiative response and $\alpha$ is the climate feedback parameter \citep[][]{andrewsForcingFeedbacksClimate2012a}.
\end{comment}

\section{Discussion}\label{sec:discussion}

% Summarize what we’ve learned
With this paper we aimed to find a common language to categorise different ways to compute model weights, describe the different choices that have to be made together with their implications, and generally highlight their advantages and disadvantages.
We differentiate between two general approaches: ensemble and individual performance weighting.
While in the former weights are computed based on the performance of a linear combination of all models, in the latter, performance is computed for each individual model. As we showed by means of toy examples, a practical advantage of the ensemble performance weighting approach is that it implicitly accounts for the issue of inter-model dependencies since each weight represents the relative contribution of the corresponding model to the weighted average instead of representing its performance as such.
We further considered two additional aspects under the ensemble performance weighting approach, namely the use of multiple diagnostics to measure performance as well as the influence of the choice of the prior distribution over weights.

% multiple diagnostics
When using multiple diagnostics, we showed that applying them jointly so that the weights are tuned to maximize the overall likelihood, i.e.~encompassing all diagnostics, yields better results in terms of the summed RMSE between the weighted average model (using the mean posterior weight vector) and the observed data than applying the diagnostics sequentially and then taking the average of the computed weight vectors. The choice of which diagnostics are most appropriate for measuring performance is a big question on its own \citep[e.g., see][]{glecklerPerformanceMetricsClimate2008} that we have not addressed here. 
Since, in the case of climate projections, we are assessing the performance of physical models, it makes sense to analyze variables that are related to the basic dynamic processes relevant for the target \citep{katzenbergerDevelopingGuidelinesWorking2025}. In any case, it should always be made clear that there is no universal answer and that any probabilistic projections are conditional on the metrics and diagnostic climate variables that were used. Furthermore, when multiple diagnostics are considered, these should be as independent of one another as possible since otherwise we would run into the same problem that we have with co-dependent models in the individual performance weighting approach: good models (diagnostics) would be erroneously overestimated while poor models (diagnostics) would be erroneously underestimated.

Note that in our toy example for the combined diagnostics we standardized the data with respect to the standard deviation of the observations so that the RMSEs are comparable across different diagnostics, i.e.~across different units. This is not necessary as we assume different likelihoods (here different variances) for each diagnostic. When scaling the data by a constant, one should be aware that the variance of the data is scaled by the square of that constant, which also decreases the variance in the weighted averages. This, in turn, decreases the differences between the likelihoods of the observed data under the weighted average models, which, thus, increases the influence of the prior. The predictions of two weighted average models, $m_{*_1}$ and $m_{*_2}$, for example, will differ less when the data is scaled, thus, the respectively used weight vectors, $w_1$ and $w_2$, yield more similar likelihoods. Put differently, when the data is scaled the influence of the likelihood is decreased. A simple possibility to verify the influence of the data is, for instance, to check how sensitive the posterior weights are with respect to different prior distributions. 

Moreover, our example of using ECS as performance metric revealed the importance of prior predictive checks: the prior should be chosen so that the data that is \emph{a priori} generated by the weighted average models is realistic and largely covers the range of plausible data.

% Add paragraph on represented uncertainties
An aspect that we did not consider in more detail in this paper is the representation of the uncertainties. The predicted uncertainty range for our target variable (e.g., future projections of the increase of global mean temperature) based on the computed posterior weight vectors comprises the uncertainty resulting from structural differences between models that is minimized within the ensemble performance weighting approach. What we have not included, but which is possible within a Bayesian formulation, is the uncertainty of the different models themselves (which can be estimated based on multiple model runs) and the uncertainty of the observations.
To include these would be important to get a more realistic estimate of the existing uncertainties.

% Caveats to BMA - when is it not appropriate
A limitation of the ensemble performance weighting approach is that it is restricted to linear combinations of individual models. 
However, in some cases it might be more appropriate to combine models in a non-linear way.
Some models, for instance, perform much better in particular regions of the globe than in others, suggesting that the weights may be modeled to differ across grid points, or regions \citep[e.g.,][]{dasbhowmikPerformancebasedMultimodelCombination2021,thaoCombiningGlobalClimate2022}. 
As the choice of the diagnostics, whether or not it is helpful to compute weights dependent on a region or even on a grid-point level also depends on the objective. If the objective is to compute the likely range of global mean temperatures in the future, it seems more plausible to stick with one weight per model whereas weights on the level of grid cells may be useful when considering likely ranges of temperatures within a certain region only.

A key focus of this work has been to compare the different assumptions and implications of the individual  and the ensemble performance weighting approach. 
Overall, past work \citep{massoudGlobalClimateModel2019,massoudBayesianModelAveraging2020} indicate that ensemble performance weighting provides benefits over individual performance weighting. As just mentioned, there may be cases in which ensemble performance weighting is not appropriate. However, there are many cases in the literature where the final goal is to obtain good performance of the combined ensemble, i.e.~we are interested in knowing the weighted-mean prediction. In those cases, ensemble performance weighting is more straightforwardly connected to the goal and provides an estimate of the uncertainty in the weights which is not accounted for in the individual performance weighting approach. If, on the other hand, the goal is simply to rank models or model versions against each other, then the individual performance weighting approach is straightforward in its application. % and should give the same result as in the ensemble weighting approach.

% Likelihood
%While we used a Gaussian likelihood here, Note though, that BMA is not limited to Gaussian weighting, but any likelihood function can be used.

% Additional considerations possible:
%Independence of observations and gridded model values
%Note Sanderson et al. (2015) section “b. Principal component analysis”, who point to EOFs as a way to retain the right scale of distances. “The use of the EOF prefilter combines fields that are trivially correlated (i.e., adjacent grid cells) into a single mode.”
%Different sources of error and combining them (eg. using a real likelihood function for performance allows explicit representation of error terms…)
%Guidelines

% Final point: BMA probably the best choice in many cases:




\conclusions  \label{sec:conclusion}%% \conclusions[modified heading if necessary]


With this paper we aimed to find a common terminology to categorise different approaches to compute model weights that facilitates a comparison between them, and makes their underlying assumptions transparent. For that we distinguish between two general methods that we call the \emph{individual performance weighting} and the \emph{ensemble performance weighting approach}. The former comprises all methods that compute a single weight vector where each model is assigned a probability relative to its individual performance while methods belonging to the latter compute a distribution over weight vectors tuned to optimize the performance of the weighted ensemble as a whole. 
The ensemble performance weighting approach, for instance, has the advantage to implicitly account for model inter-dependencies. However, which method is more appropriate will be different from case to case depending on the target application of the model weights. When using the ensemble performance weighting approach, we showed that a careful assessment of the prior distribution is necessary to draw appropriate conclusions about the resulting posterior weights and their application to target variables. Overall, our analyses showed the importance of explicitly acknowledging the influence of our choices, starting with the general approach to compute model weights. 


%Weights based on the latter approach have the benefit that they already contain information about possible co-dependencies between the models in the ensemble which have to be accounted for separately in the individual weighting approach.
%Furthermore, this approach naturally incorporates uncertainties of the weights instead of assuming a single ``true'' weight vector.
%We therefore highly recommend using the ensemble performance weighting approach to compute model weights.
%We have shown that, under this approach, it is further advisable, when using multiple diagnostics, to apply them jointly instead of sequentially and then averaging the mean weights.




%% The following commands are for the statements about the availability of data sets and/or software code corresponding to the manuscript.
%% It is strongly recommended to make use of these sections in case data sets and/or software code have been part of your research the article is based on.

%\codeavailability{TEXT} %% use this section when having only software code available


%\dataavailability{TEXT} %% use this section when having only data sets available


\codedataavailability{All code and data used in the paper is publicly available under \url{https://github.com/awi-esc/WeightedEnsembles}. The CMIP data was downloaded from ESGF and preprocessed (regridded) with ESMValTool. }%The ESMValTool recipes are provided in the GitHub repository.} %% use this section when having data sets and software code available


%\sampleavailability{TEXT} %% use this section when having geoscientific samples available

%\videosupplement{TEXT} %% use this section when having video supplements available


\authorcontribution{B.G. and A.R. developed the initial concept for the paper. B.G. performed the analysis, prepared the figures and wrote the manuscript, with contributions from all authors.} %% this section is mandatory

\competinginterests{The authors declare that they have no conflict of interest.} %% this section is mandatory even if you declare that no competing interests are present

%\disclaimer{TEXT} %% optional section

\begin{acknowledgements}
B.G., M.P. and A.R. received funding from the European Union (ERC, FORCLIMA, 101044247).
\end{acknowledgements}


\appendix \label{sec:appendix}
%\section{Bayesian Model Averaging} \label{sec:appendix-bma}   %% Appendix A
%TODO: Describe in statistical literature and how it differs from how it is normallly used with climate models.

\section{Influence of prior on weighted projections}

Figure \ref{fig:appendix-projections-mmm-bma} shows the influence of using a prior distribution over weight vectors in the ensemble performance weighting approach that strongly restricts the considered data. Here, we use a Dirichlet prior with concentration parameter $\alpha=1$ that yields weight vectors with roughly equal weights assigned to every model. This strongly restricts the predictions of the weighted average models that are considered \emph{a priori} as shown in Fig.~\ref{fig:expected-ecs-dirichlet}b in the main text. This, in turn, is reflected in Fig.~\ref{fig:appendix-projections-mmm-bma} where the shown inner 90\% uncertainty range for the projected increase in global mean temperature is nearly identical when using the prior weight vectors as compared to using the posterior weight vectors. Therefore, this extremely narrow range is not due to the data: we get almost the same uncertainty range simply by using the chosen prior weights without any consideration of the likelihood function.

\begin{figure}[H]
    \centering
    %\includegraphics{sections/figures/appendix/projections-bma-mmm-with-prior-1.pdf}
    \includegraphics{sections/figures/figA1.pdf}
    %\caption{Predicted increase in global mean near-surface air temperature with respect to the time period from 1995--2014 until the end of the century, under emission scenario SSP5.85. Solid lines show the multi-model mean (brown) of an ensemble of $N=38$ CMIP6 models and the weighted average of this ensemble (blue) using the mean posterior weights applying the ensemble performance weighting approach with prior Dirichlet($\mathbf{1}$) and with the target distribution from Fig.~\ref{fig:expected-ecs-dirichlet} as likelihood. The green dashed lines show the weighted mean (thick line), using the mean prior weights. Uncertainty bands (for the prior dashed green lines) represent the 0.05 and 0.95 (weighted) quantiles.}
    \caption{Like Fig.~\ref{fig:projections-mmm-bma} in the main text, but when using a Dirichlet prior with $\alpha=1$.}
    \label{fig:appendix-projections-mmm-bma}
\end{figure}

%\subsection{}     %% Appendix A1, A2, etc.

%\footnote{Note that similarly to the case when $\sigma^2$ is sampled in the individual performance weighting approach, when using multiple diagnostics, it may happen that $\sigma_d^2$ for a diagnostic $d$ receives large values so that the observed data for this diagnostic is nearly equally likely under all weighted average models and all  }0

\section{}\label{sec:appendix-B}

Here we show examples of different weighting approaches that yield the same ranking order and only differ with respect to the magnitude of the individual weights. Note that for simplicity, we left out the area-weights here that should be applied for spatial data.

\begin{enumerate}
    \item Assume i.i.d Gaussian errors, i.e. $y(s) = m_i(s) + \epsilon$ with $\epsilon \sim \mathcal{N}(0,\sigma^2)$ and weights for model $i$ are proportional to the likelihood of the data under model $i$
    \begin{itemize}
        \item That is, we assume that each data point (y(s)) comes from a normal distribution with the prediction of model $i$ ($m_i$) as mean and with standard deviation $\sigma$.
        \item The likelihood of the observed data under model $i$ is then given by $P(Y\mid M_i, \sigma) = \prod_{s} {\frac {1}{\sqrt {2\pi \sigma ^{2}}}}\exp {\left(-{\frac {(y(s) - m_i(s)^{2}}{2\sigma ^{2}}}\right)}$, iterating over all data points, e.g. all spatial locations $s$.
        \item The log likelihood of all observed data is then given by: $-{\frac {n}{2}}\log(2\pi \sigma ^{2})-{\frac {1}{2\sigma ^{2}}}\sum _{i=1}^{n}\left(y(s)-m_i(s) \right)^{2}$
        \item turned into a likelihood again: $w_i \propto \exp(c_1) - \exp(\frac{1}{2\sigma^2} \cdot \sum _{i=1}^{n}\left(y(s)-m_i(s) \right)^{2} = \exp(c_1) - \exp(c_2 \cdot \textrm{SSE})$; i.e. $w_i\propto \exp(-c_2\cdot SSE)$ where SSE is the sum of squared errors and $c_1, c_2$ are constants.
    \end{itemize}

\item \citet{sierraInformationtheoreticApproachObtain2025}: $w_i \propto \frac{1}{\hat{\sigma}^2}$ with $\hat{\sigma}^2 = \frac{\sum_{s} (m_i(s) - y(s))^2}{n} = \frac{\textrm{SSE}}{n} = \textrm{MSE}$. Therefore, $w_i \propto n \cdot \frac{1}{\textrm{SSE}}=\frac{1}{\textrm{MSE}}$ with MSE referring to the mean squared error.

\item \citet{brunnerReducedGlobalWarming2020} performance weights for a single diagnostic variable: $w_i \propto \exp(-(\frac{\sqrt{\textrm{MSE}}}{\sigma_D})^2)$ where $\sigma_D$ is a tuned hyperparameter. 

\end{enumerate}


\begin{comment}
\section{Bayesian Model Averaging}\label{sec:appendix-bma}
In the introduction we mentioned that some papers have referred to what we call the ‘Ensemble performance weighting’ approach as ‘Bayesian Model Averaging’, short BMA. However, in the statistical literature, BMA has a different meaning that we believe is worth disentangling and setting in relation to the ensemble performance weighting approach that we have described here. In the classical sense, BMA applies Bayesian inference on the problem of model selection by taking into account the uncertainties about the models themselves instead of assuming one specific model \citep[for a review on BMA see][]{fragosoBayesianModelAveraging2018}. For a concrete statistical model $\mathcal{M}$ with certain parameters, summarized as $\theta_{\mathcal{M}}$, we have the following quantities:

\begin{itemize}
    \item[\textbullet] Posterior distribution: $P(\theta_{\mathcal{M}} \mid Y)$
    \item[\textbullet] Likelihood: $P(Y\mid \theta_{\mathcal{M}})$
    \item[\textbullet] Posterior Predictive Distribution for new data $\tilde{y}$:
        \begin{equation}
            P(\tilde{y}\mid Y, \mathcal{M}) = \int_{\theta_{\mathcal{M}}} P(\hat{y}\mid \theta_{\mathcal{M}}) P(\theta_{\mathcal{M}}\mid Y) d\theta_{\mathcal{M}}
        \end{equation}
\end{itemize}

However, we may not be certain that this model is the best representation for our data. BMA now takes into account this uncertainty about the models themselves when making predictions about new data $\tilde{y}$: 

\begin{equation}
    P(\tilde{y}\mid Y) = \sum_{i=1}^N P(\tilde{y}\mid Y, \mathcal{M}_i) \cdot P(\mathcal{M}_i \mid Y) 
\end{equation}

where the weights, $P(\mathcal{M}_i\mid Y)$, are computed via the marginal likelihood:

\begin{equation}
    P(\mathcal{M}_i\mid Y) \propto P(Y\mid \mathcal{M}_i) = \int_{\theta_{\mathcal{M}}} P(Y\mid \theta_{\mathcal{M},i}) \cdot P(\theta_{\mathcal{M},i}) d\theta_{\mathcal{M}} 
\end{equation}

This, in turn, means that in classical BMA, models are interpreted as representing mutually exclusive hypotheses, including the `correct' (or `most correct') model.  In the long run, after collecting more and more data, all probability mass will be put on the best model. This clashes with the interpretation of an ensemble of climate models which all contribute to our understanding of the climate system and among which we do not believe to find a single `correct' model. Instead, climate models are considered to be inter-dependent imperfect representations of the climate system.
Contrary to the `Ensemble Performance Weighting' approach that we have described, in the classical definition of BMA, the weights are posterior model weights, i.e.~they refer to individual models.

\citet{rafteryUsingBayesianModel2005} extended the statistical definition of BMA to dynamical models, more precisely to ensembles of weather forecast models. They combine the probability densities of the models' forecasts into a mixture distribution via mixture weights that are computed with the Expectation-Maximization (EM) algorithm. The weights are interpreted as the respective forecast to be the best forecast in the ensemble. That is, the 
objective in this 'BMA-approach' is to maximize the joint mixture distribution, however, assuming that each observation originates from one of the component distributions. That is, for every observation, a latent variable is assumed that specifies from which individual model it was generated. The computed weight vector thus summarizes how often each model is the best among all models in the ensemble, in explaining the observed data. 
This contrasts the `Ensemble Performance Weighting' approach which is based on the assumption that a weighted average model generated the observed data and which computes a distribution over weight vectors, instead of estimating a single weight vector. 

There are applications following the classical as well as the approach from \citet{rafteryUsingBayesianModel2005} to BMA in the context climate models \citep[e.g.,][]{zhuFutureProjectionsUncertainty2013,olsonSimpleMethodBayesian2016, fanBayesianPosteriorPredictive2017, zhuIntercomparisonMultimodelEnsembleprocessing2023}.


%\begin{equation}
 %   P(\tilde{y}\mid Y) = \sum_{i=1}^N w_i \cdot P(\tilde{y}\mid \mathcal{M}_i)
%\end{equation}


%Contrary to climate projections, weather forecasts work on very short timescales, thus

\end{comment}


\noappendix       %% use this to mark the end of the appendix section. Otherwise the figures might be numbered incorrectly (e.g. 10 instead of 1).

%% Regarding figures and tables in appendices, the following two options are possible depending on your general handling of figures and tables in the manuscript environment:

%% Option 1: If you sorted all figures and tables into the sections of the text, please also sort the appendix figures and appendix tables into the respective appendix sections.
%% They will be correctly named automatically.

%% Option 2: If you put all figures after the reference list, please insert appendix tables and figures after the normal tables and figures.
%% To rename them correctly to A1, A2, etc., please add the following commands in front of them:

\appendixfigures  %% needs to be added in front of appendix figures



\appendixtables   %% needs to be added in front of appendix tables

%% Please add \clearpage between each table and/or figure. Further guidelines on figures and tables can be found below.




%% REFERENCES

%% Since the Copernicus LaTeX package includes the BibTeX style file copernicus.bst,
%% authors experienced with BibTeX only have to include the following two lines:
%%
\bibliographystyle{copernicus}
\bibliography{model-weighting-paper.bib}

%%
%% URLs and DOIs can be entered in your BibTeX file as:
%%
%% URL = {http://www.xyz.org/~jones/idx_g.htm}
%% DOI = {10.5194/xyz}


%% LITERATURE CITATIONS
%%
%% command                        & example result
%% \citet{jones90}|               & Jones et al. (1990)
%% \citep{jones90}|               & (Jones et al., 1990)
%% \citep{jones90,jones93}|       & (Jones et al., 1990, 1993)
%% \citep[p.~32]{jones90}|        & (Jones et al., 1990, p.~32)
%% \citep[e.g.,][]{jones90}|      & (e.g., Jones et al., 1990)
%% \citep[e.g.,][p.~32]{jones90}| & (e.g., Jones et al., 1990, p.~32)
%% \citeauthor{jones90}|          & Jones et al.
%% \citeyear{jones90}|            & 1990



%% FIGURES

%% When figures and tables are placed at the end of the MS (article in one-column style), please add \clearpage
%% between bibliography and first table and/or figure as well as between each table and/or figure.

% The figure files should be labelled correctly with Arabic numerals (e.g. fig01.jpg, fig02.png).


%% ONE-COLUMN FIGURES

%%f
%\begin{figure}[t]
%\includegraphics[width=8.3cm]{FILE NAME}
%\caption{TEXT}
%\end{figure}
%
%%% TWO-COLUMN FIGURES
%
%%f
%\begin{figure*}[t]
%\includegraphics[width=12cm]{FILE NAME}
%\caption{TEXT}
%\end{figure*}
%
%
%%% TABLES
%%%
%%% The different columns must be seperated with a & command and should
%%% end with \\ to identify the column brake.
%
%%% ONE-COLUMN TABLE
%
%%t
%\begin{table}[t]
%\caption{TEXT}
%\begin{tabular}{column = lcr}
%\tophline
%
%\middlehline
%
%\bottomhline
%\end{tabular}
%\belowtable{} % Table Footnotes
%\end{table}
%
%%% TWO-COLUMN TABLE
%
%%t
%\begin{table*}[t]
%\caption{TEXT}
%\begin{tabular}{column = lcr}
%\tophline
%
%\middlehline
%
%\bottomhline
%\end{tabular}
%\belowtable{} % Table Footnotes
%\end{table*}
%
%%% LANDSCAPE TABLE
%
%%t
%\begin{sidewaystable*}[t]
%\caption{TEXT}
%\begin{tabular}{column = lcr}
%\tophline
%
%\middlehline
%
%\bottomhline
%\end{tabular}
%\belowtable{} % Table Footnotes
%\end{sidewaystable*}
%
%
%%% MATHEMATICAL EXPRESSIONS
%
%%% All papers typeset by Copernicus Publications follow the math typesetting regulations
%%% given by the IUPAC Green Book (IUPAC: Quantities, Units and Symbols in Physical Chemistry,
%%% 2nd Edn., Blackwell Science, available at: http://old.iupac.org/publications/books/gbook/green_book_2ed.pdf, 1993).
%%%
%%% Physical quantities/variables are typeset in italic font (t for time, T for Temperature)
%%% Indices which are not defined are typeset in italic font (x, y, z, a, b, c)
%%% Items/objects which are defined are typeset in roman font (Car A, Car B)
%%% Descriptions/specifications which are defined by itself are typeset in roman font (abs, rel, ref, tot, net, ice)
%%% Abbreviations from 2 letters are typeset in roman font (RH, LAI)
%%% Vectors are identified in bold italic font using \vec{x}
%%% Matrices are identified in bold roman font
%%% Multiplication signs are typeset using the LaTeX commands \times (for vector products, grids, and exponential notations) or \cdot
%%% The character * should not be applied as mutliplication sign
%
%
%%% EQUATIONS
%
%%% Single-row equation
%
%\begin{equation}
%
%\end{equation}
%
%%% Multiline equation
%
%\begin{align}
%& 3 + 5 = 8\\
%& 3 + 5 = 8\\
%& 3 + 5 = 8
%\end{align}
%
%
%%% MATRICES
%
%\begin{matrix}
%x & y & z\\
%x & y & z\\
%x & y & z\\
%\end{matrix}
%
%
%%% ALGORITHM
%
%\begin{algorithm}
%\caption{...}
%\label{a1}
%\begin{algorithmic}
%...
%\end{algorithmic}
%\end{algorithm}
%
%
%%% CHEMICAL FORMULAS AND REACTIONS
%
%%% For formulas embedded in the text, please use \chem{}
%
%%% The reaction environment creates labels including the letter R, i.e. (R1), (R2), etc.
%
%\begin{reaction}
%%% \rightarrow should be used for normal (one-way) chemical reactions
%%% \rightleftharpoons should be used for equilibria
%%% \leftrightarrow should be used for resonance structures
%\end{reaction}
%
%
%%% PHYSICAL UNITS
%%%
%%% Please use \unit{} and apply the exponential notation (e.g. 20\,\unit{W\,m^{-2}})


\end{document}
